% Independently audited and accepted by the hardened five-source Rethlas v3.3.3-r1 path.
% Exact mathematical payload: run 20260901T070930Z_773fc66f34a8_35b1bb24.
\providecommand{\an}{\mathrm{an}}

\begin{proposition}\label{prop:p053-projection-recursion-repaired-candidate}
Let \(d\geq1\), and let \((\mathbf X,[\mathbf X])\) be a pure
\(d\)-dimensional balanced rational polyhedral space, endowed with its
boundaryless tropical-space structure and its integration current
\(\delta_{[\mathbf X]}\).  Let \(Y\) be a pure \(d\)-dimensional
boundaryless CLD tropical space, let \([Y]\) be its positive closed current
of integration, and let

\[
 h:Y\longrightarrow\mathbf X
\]

be a topologically proper morphism of tropical spaces.  Assume that

\begin{equation}
 h_*[Y]=\delta_{[\mathbf X]}
 \label{eq:p053-repaired-top-calibration}
\end{equation}

and that

\begin{equation}
 h^*\operatorname{PAPC}(\mathbf X)
 \subseteq \operatorname{CPsh}(Y).
 \label{eq:p053-repaired-psh-pullback}
\end{equation}

For \(a_1,\ldots,a_d\in\operatorname{PAPC}(\mathbf X)\), set

\begin{equation}
 S_0=[Y],\qquad
 S_j=\dprimeddprime\bigl((a_j\circ h)S_{j-1}\bigr)
 \quad (1\leq j\leq d),
 \label{eq:p053-repaired-source-recursion}
\end{equation}

and

\begin{equation}
 c_0=[\mathbf X],\qquad c_j=a_j\cdot c_{j-1}.
 \label{eq:p053-repaired-target-recursion}
\end{equation}

Then every \(S_j\) is positive and closed, every \(c_j\) is an effective
balanced tropical cycle, and

\begin{equation}
 h_*S_j=\delta_{c_j}
 \quad (0\leq j\leq d).
 \label{eq:p053-repaired-intermediate-pushforward}
\end{equation}

In particular,

\begin{equation}
 h_*\!\left(
   \dprimeddprime[a_1\circ h]\wedge\cdots\wedge
   \dprimeddprime[a_d\circ h]\wedge[Y]
 \right)
 =\delta_{a_1\cdots a_d\cdot[\mathbf X]}
 =\MA_{\mathrm{poly}}(a_1,\ldots,a_d).
 \label{eq:p053-repaired-pa-comparison}
\end{equation}

Both iterated products in \eqref{eq:p053-repaired-pa-comparison} are
independent of the displayed order.

Let now \(\gamma\) be a piecewise-affine polyhedrally convex reference
function that is positive at infinity.  If

\[
 a_1,\ldots,a_d\in\mathrm{PC}_{\mathrm{reg}}(\mathbf X,\gamma),
\]

then every \(a_i\circ h\) is continuous psh and

\begin{equation}
 h_*\!\left(
   \dprimeddprime[a_1\circ h]\wedge\cdots\wedge
   \dprimeddprime[a_d\circ h]\wedge[Y]
 \right)
 =\MA_{\mathrm{poly}}(a_1,\ldots,a_d)
 \label{eq:p053-repaired-regularizable-comparison}
\end{equation}

as Radon measures on \(\mathbf X\).  For every simplicial polyhedral
structure \(\Pi\) carrying \(\gamma\), the measure on the right, viewed on
the associated BMPS compactification \(\mathbf X_\Pi\), has zero mass on
\(\mathbf X_\Pi\setminus\mathbf X\).
\end{proposition}

\begin{corollary}\label{cor:p053-algebraic-chart-repaired-candidate}
Let \(K\) be an algebraically closed complete field with a nontrivial
non-Archimedean absolute value, and let \(U\) be an irreducible very affine
\(K\)-variety of dimension \(d\).  Use the convention

\[
 h=\operatorname{trop}_U=(-\log|f_1|,\ldots,-\log|f_n|):
 U^{\an}\longrightarrow\mathbf X:=\operatorname{Trop}(U)
\]

for a closed embedding of \(U\) in a split torus.  On the CLD side this is
the sign-reversed moment chart.  Equip \(\mathbf X\) with its saturated
lattices and analytic tropical multiplicities, and assume that it is a
rational balanced polyhedral space in the sense of
Botero--Mazzon--Pille-Schneider.  Then
\eqref{eq:p053-repaired-top-calibration} and
\eqref{eq:p053-repaired-psh-pullback} hold for \(Y=U^{\an}\).  Consequently,
all conclusions of
\Cref{prop:p053-projection-recursion-repaired-candidate} apply.

For \(a_i\in\operatorname{PAPC}(\mathbf X)\), put \(u_i=a_i\circ h\),
endow the trivial line bundle with the piecewise-smooth delta-metric

\[
 \lVert1\rVert_i=\mathrm e^{-u_i},
\]

and write \(\Theta_i=c_1(\mathcal O_U,\lVert\cdot\rVert_i)\) for its first
Chern delta-form.  Each \(\Theta_i\) is a positive delta-form.  More
precisely, for every ordering and every \(1\leq j\leq d\), the product of
the first \(j\) factors is a positive generalized delta-form; this assertion
is made for these particular Chern forms and not for arbitrary products of
positive generalized delta-forms.  The top-degree delta-form
\(\Theta_1\wedge\cdots\wedge\Theta_d\) therefore induces a positive,
locally finite Radon measure \(\mu_\delta\) on \(U^{\an}\), and

\begin{equation}
 [\mu_\delta]
 =\dprimeddprime[u_1]\wedge\cdots\wedge
   \dprimeddprime[u_d]\wedge[U^{\an}],
 \qquad
 h_*\mu_\delta=\MA_{\mathrm{poly}}(a_1,\ldots,a_d).
 \label{eq:p054-repaired-delta-comparison}
\end{equation}

No finiteness of the total mass of \(\mu_\delta\) is asserted without an
additional properness or support hypothesis.

More generally, take
\(a_i\in\mathrm{PC}_{\mathrm{reg}}(\mathbf X,\gamma)\) and arbitrary
decreasing approximations

\[
 a_{i,k}\in\operatorname{PAPC}(\mathbf X,\gamma),
 \qquad a_{i,k}\searrow a_i.
\]

Let \(\mu_{\delta,k}\) be the locally finite delta-form measure associated
with \(a_{1,k},\ldots,a_{d,k}\).  Then \(\mu_{\delta,k}\) converges vaguely
on \(U^{\an}\) to

\begin{equation}
 \mu_{\mathrm{CLD}}
 :=\dprimeddprime[u_1]\wedge\cdots\wedge
   \dprimeddprime[u_d]\wedge[U^{\an}].
 \label{eq:p054-repaired-source-limit}
\end{equation}

The limit is independent of all approximating sequences and satisfies

\begin{equation}
 h_*\mu_{\mathrm{CLD}}
 =\MA_{\mathrm{poly}}(a_1,\ldots,a_d).
 \label{eq:p054-repaired-limit-pushforward}
\end{equation}

No assertion is made that \(\mu_{\mathrm{CLD}}\) is represented by a finite
delta-form or has finite total mass without an additional properness or
support hypothesis.
\end{corollary}

\begin{proof}[of
\Cref{prop:p053-projection-recursion-repaired-candidate}]
We first compare the two target-side divisor conventions.  Let \(c\) be a
positive balanced tropical cycle and let \(a\) be piecewise affine on a
common rational subdivision.  If \(\tau\) is a codimension-one face, both
\cite[Definition~1.10]{GK17} and \cite[Definition~2.31]{BMPS26} assign the
weight

\begin{equation}
 (a\cdot c)(\tau)
 =\sum_{\sigma\succ\tau}c(\sigma)
       a_\sigma(n_{\sigma/\tau})
  -a_\tau\!\left(
       \sum_{\sigma\succ\tau}c(\sigma)n_{\sigma/\tau}
    \right).
 \label{eq:p053-repaired-weight-match}
\end{equation}

Here \(n_{\sigma/\tau}\) is the primitive outward class in the saturated
quotient lattice.  Thus the two conventions have the same sign and use the
same primitive saturated-quotient normalization.  Lattice indices enter
tropical pushforwards, not the displayed divisor convention.

The weighted piecewise-linear Poincar\'e--Lelong formula
\cite[Corollary~3.20]{GK17}, applied after the chosen subdivision, gives

\begin{equation}
 \dprimeddprime(a\delta_c)=\delta_{a\cdot c}.
 \label{eq:p053-repaired-weighted-pl}
\end{equation}

Indeed, against a compactly supported smooth test form \(\eta\), the two
sides of \eqref{eq:p053-repaired-weighted-pl} have values

\[
 \int_c a\,\dprimeddprime\eta
 \quad\text{and}\quad
 \int_{a\cdot c}\eta,
\]

respectively, and these are equal by the cited formula.  The facewise smooth
Hessian term in \cite[Theorem~3.16]{GK17} vanishes because \(a\) is affine on
each cell.  Formula \eqref{eq:p053-repaired-weighted-pl} holds for every
balanced intermediate cycle, not only for the fundamental cycle.

We next establish the projection identity used in the induction.  Let \(S\)
be a positive current of bidegree \((p,p)\) on \(Y\), with \(p<d\), and let
\(b\) be a continuous function on \(\mathbf X\).  A positive current has
order zero, so both \((b\circ h)S\) and \(b h_*S\) are defined.  For a
compactly supported smooth test form \(\eta\) of bidegree
\((d-p-1,d-p-1)\), properness of \(h\), functoriality of pullback, and the
definition of the current differentials give

\begin{align}
 \left\langle h_*\dprimeddprime((b\circ h)S),\eta\right\rangle
 &=\left\langle (b\circ h)S,
                   \dprimeddprime(h^*\eta)\right\rangle \\
 &=\left\langle S,h^*(b\,\dprimeddprime\eta)\right\rangle \\
 &=\left\langle \dprimeddprime(bh_*S),\eta\right\rangle.
 \label{eq:p053-repaired-continuous-projection}
\end{align}

There is no sign in this identity: two first-order current differentials are
applied.  Hence

\begin{equation}
 h_*\dprimeddprime((b\circ h)S)
 =\dprimeddprime(bh_*S).
 \label{eq:p053-repaired-projection-identity}
\end{equation}

This is the computation in the proof of
\cite[Proposition~8.5.6]{CLD25}; continuity is enough for the equality, while
the psh assumption in that proposition is what supplies positivity.  Compact
support and topological properness justify every pullback of a test form.
Because the source and target both have pure dimension \(d\), pushforward
preserves the displayed bidegrees.  Boundarylessness and compact support
exclude an additional boundary term.

We prove \eqref{eq:p053-repaired-intermediate-pushforward} by induction.  It
holds for \(j=0\) by \eqref{eq:p053-repaired-top-calibration}.  Suppose it
holds for \(j-1\).  By \eqref{eq:p053-repaired-psh-pullback} and
\cite[Lemma~8.5.5]{CLD25}, the current

\[
 S_j=\dprimeddprime((a_j\circ h)S_{j-1})
\]

is positive and closed.  Since \(c_{j-1}\) is positive, the defining PAPC
condition says that \(c_j=a_j\cdot c_{j-1}\) is effective; it is balanced by
the tropical divisor calculus.  Equations
\eqref{eq:p053-repaired-projection-identity} and
\eqref{eq:p053-repaired-weighted-pl} now yield

\begin{align}
 h_*S_j
 &=\dprimeddprime(a_jh_*S_{j-1}) \\
 &=\dprimeddprime(a_j\delta_{c_{j-1}}) \\
 &=\delta_{a_j\cdot c_{j-1}}
 =\delta_{c_j}.
 \label{eq:p053-repaired-projection-induction}
\end{align}

Here \(S_j\) and \(\delta_{c_j}\) have bidegree \((j,j)\), while \(c_j\)
has pure dimension \(d-j\).  No current depending only on the abstract
target cycle \(c_{j-1}\) is being constructed on \(Y\); the induction uses
the actual source current \(S_{j-1}\).

For \(j=d\), the cycle \(c_d\) is zero-dimensional, and
\cite[Definition~4.7]{BMPS26} identifies it, with its displayed weights and
without a factorial, with
\(\MA_{\mathrm{poly}}(a_1,\ldots,a_d)\).  This proves
\eqref{eq:p053-repaired-pa-comparison}.  Commutativity of the target
divisors is \cite[Proposition~1.14]{GK17}, equivalently
\cite[Lemma~2.35]{BMPS26}; symmetry of the source product is
\cite[Proposition~8.7.4]{CLD25}.

We turn to regularizable weights.  For every \(i\), choose

\[
 a_{i,k}\in\operatorname{PAPC}(\mathbf X,\gamma),
 \qquad a_{i,k}\searrow a_i,
\]

uniformly as in \cite[Definition~3.23]{BMPS26}.  Their pullbacks decrease
locally uniformly on \(Y\).  The local uniform closure of continuous psh
functions shows that \(a_i\circ h\) is continuous psh.  The joint continuity
of mixed CLD products, \cite[Proposition~8.7.2]{CLD25}, gives

\begin{equation}
 \bigwedge_{i=1}^d\dprimeddprime[a_{i,k}\circ h]\wedge[Y]
 \longrightarrow
 \bigwedge_{i=1}^d\dprimeddprime[a_i\circ h]\wedge[Y]
 \label{eq:p053-repaired-cld-limit}
\end{equation}

as currents.  On the target, \cite[Theorem~4.13]{BMPS26} gives

\begin{equation}
 \MA_{\mathrm{poly}}(a_{1,k},\ldots,a_{d,k})
 \longrightarrow
 \MA_{\mathrm{poly}}(a_1,\ldots,a_d).
 \label{eq:p053-repaired-poly-limit}
\end{equation}

Let \(\eta\) be a compactly supported smooth target test function.
Properness makes \(h^*\eta\) compactly supported.  We may therefore pass to
the limit in the piecewise-affine equality and obtain
\eqref{eq:p053-repaired-regularizable-comparison} on smooth tests.  Smooth
compactly supported functions are dense in continuous compactly supported
functions by \cite[Theorem~6.2.5 and Corollary~6.2.6]{CLD25}; since both sides
are positive Radon measures, equality extends to all continuous compact
tests.

It remains to justify the asserted mixed boundary behavior.  Put

\[
 \bar a=\frac1d\sum_{i=1}^d a_i.
\]

The class \(\mathrm{PC}_{\mathrm{reg}}(\mathbf X,\gamma)\) is convex, so
\(\bar a\) lies in it.  By \cite[Proposition~4.16]{BMPS26}, the diagonal
measure \(\MA_{\mathrm{poly}}(\bar a,\ldots,\bar a)\), viewed on any
\(\mathbf X_\Pi\) occurring in the statement, gives no mass to
\(\mathbf X_\Pi\setminus\mathbf X\).  Multilinearity and positivity give

\begin{equation}
 \MA_{\mathrm{poly}}(\bar a,\ldots,\bar a)
 =d^{-d}\!\sum_{1\leq i_1,\ldots,i_d\leq d}
   \MA_{\mathrm{poly}}(a_{i_1},\ldots,a_{i_d}).
 \label{eq:p053-repaired-mixed-boundary-expansion}
\end{equation}

Every summand is a positive measure.  Since their sum has zero boundary
mass, each summand has zero boundary mass, in particular the summand with
arguments \(a_1,\ldots,a_d\).  This completes the proof.
\end{proof}

\begin{proof}[of
\Cref{cor:p053-algebraic-chart-repaired-candidate}]
The tropicalization map \(h\) is topologically proper.  This is precisely
the properness input used in the proof of \cite[Proposition~4.21]{GK17},
where it is cited from BPR16, Remark~2.3.  If \(\alpha\) is a compactly
supported smooth form of top bidegree on \(\mathbf X\), the tropical
integration formula \cite[Definition~5.8]{GK17} gives

\[
 \left\langle h_*[U^{\an}],\alpha\right\rangle
 =\int_{U^{\an}}h^*\alpha
 =\int_{\operatorname{Trop}(U)}\alpha
 =\left\langle\delta_{[\mathbf X]},\alpha\right\rangle.
\]

The analytic tropical multiplicities already contain the local degrees, so
no further factor occurs.  The sign change from the basic CLD convention
\((\log|f_i|)_i\) to \(( -\log|f_i|)_i\) is the integral-linear map
\(-\operatorname{id}\).  It has lattice index one, and its action on a top
\((d,d)\)-form has sign \((-1)^{2d}=1\).  Thus
\eqref{eq:p053-repaired-top-calibration} holds with the stated convention.

We next prove \eqref{eq:p053-repaired-psh-pullback}, including the
functorial form needed below.  Fix
\(a\in\operatorname{PAPC}(\mathbf X)\).  We first record the local
calculation that will be used on every compatible tropical frame.  Let

\[
 F:N'_{\mathbf R}\longrightarrow N_{\mathbf R}
\]

be an integral-affine map, let \(C'\) be a rational polyhedral complex in
\(N'_{\mathbf R}\) mapped into \(\mathbf X\), and let \(z\) be any positive
balanced \(k\)-cycle on an open subset of \(C'\).  Refine the source and
target complexes so that \(F\) is cellular and \(a\) and
\(\varphi:=a\circ F\) are affine on their cells.  We claim that
\(\varphi\cdot z\) is effective.

Let \(\tau'\) be a \((k-1)\)-cell of \(z\), and choose
\(\xi\in\operatorname{relint}(\tau')\).  A further common refinement lets
us assume that \(y:=F(\xi)\) lies in the relative interior of a target cell
\(\delta\), that \(F(\tau')\subseteq\delta\), and that every cell adjacent
to \(\tau'\) maps into a target cell containing \(\delta\).  Put

\[
 N_{\tau'}=N'\cap\operatorname{span}_{\mathbf R}(\tau'-\tau'),
 \qquad
 N_\delta=N\cap\operatorname{span}_{\mathbf R}(\delta-\delta).
\]

These are saturated lattices.  Since \(F\) is cellular, its linear part
sends \(N_{\tau'}\) into \(N_\delta\) and therefore induces an integral
linear map of torsion-free quotient lattices

\begin{equation}
 \overline F:N'/N_{\tau'}\longrightarrow N/N_\delta.
 \label{eq:p053-vthree-quotient-map}
\end{equation}

Modulo the lineality space of \(\tau'\), the local star of \(z\) is a
positive balanced one-cycle \(z_{\tau'}\).  Its ray corresponding to a
\(k\)-cell \(\sigma'\succ\tau'\) has weight
\(m_{\sigma'}:=z(\sigma')\) and primitive generator
\(\bar n_{\sigma'/\tau'}\).  We write
\(N_{\sigma'/\tau'}:=N_{\sigma'}/N_{\tau'}\) for the resulting
rank-one ray lattice; it is generated by
\(\bar n_{\sigma'/\tau'}\).  The local conical germ of \(a\) at \(y\) is
given on a cone \(C_{\sigma/\delta}\) by \(a_\sigma-a_\delta\), as in
\cite[Definition~2.30]{BMPS26}.  It vanishes on \(N_{\delta,\mathbf R}\)
and hence descends to a piecewise-linear function \(\bar a_y\) on the
target quotient star.  The source germ is exactly
\(\bar a_y\circ\overline F\).

For every noncontracted source ray, write

\begin{equation}
 \overline F(\bar n_{\sigma'/\tau'})
 =e_{\sigma'}n_{\nu(\sigma')},
 \qquad
 e_{\sigma'}
 =[N_{\nu(\sigma')}:
   \overline F(N_{\sigma'/\tau'})]\in\mathbf Z_{>0},
 \label{eq:p053-vthree-ray-index}
\end{equation}

where \(n_{\nu(\sigma')}\) is the primitive generator of the image ray.
Contracted rays have zero image and are omitted from this notation.  The
pushforward convention \cite[(1.4.3)]{GK17} assigns the image ray \(\nu\)
the weight

\begin{equation}
 m_\nu
 =\sum_{\nu(\sigma')=\nu}e_{\sigma'}m_{\sigma'}.
 \label{eq:p053-vthree-pushforward-weight}
\end{equation}

Thus all multiplicities and lattice indices are nonnegative and are retained
in the pushforward.  Formula \eqref{eq:p053-repaired-weight-match}, applied
in the saturated quotient lattices, now gives the exact coefficient identity

\begin{align}
 (\varphi\cdot z)(\tau')
 &=\sum_{\sigma'\ \mathrm{noncontracted}}
      m_{\sigma'}\,
      \bar a_y\!\left(
         \overline F(\bar n_{\sigma'/\tau'})
      \right) \\
 &=\sum_\nu
      \left(\sum_{\nu(\sigma')=\nu}
         e_{\sigma'}m_{\sigma'}\right)\bar a_y(n_\nu) \\
 &=\bigl(\bar a_y\cdot\overline F_*z_{\tau'}\bigr)(0).
 \label{eq:p053-vthree-corner-coefficient}
\end{align}

The first equality includes no hidden tangential term: passing from \(a\) to
its conical germ subtracts the affine restriction on \(\delta\), and balancing
annihilates its pullback.  A contracted ray contributes zero because its
primitive source generator maps to zero.  Equivalently,
\eqref{eq:p053-vthree-corner-coefficient} is the tropical divisor projection
formula

\begin{equation}
 \overline F_*\bigl(
   (\bar a_y\circ\overline F)\cdot z_{\tau'}
 \bigr)
 =\bar a_y\cdot\overline F_*z_{\tau'},
 \label{eq:p053-vthree-local-projection}
\end{equation}

from \cite[Proposition~1.13]{GK17}.  Pushforward of the origin appearing on
the left has index \([0:0]=1\); this is why the two coefficients, rather than
only their positive multiples, are equal.

If every ray is contracted, \eqref{eq:p053-vthree-corner-coefficient} is
zero.  Otherwise \(\overline F_*z_{\tau'}\) is a positive balanced
one-cycle on the quotient target star.  After refining so that its rays are
cells, it lifts to a positive local cycle of dimension
\(\dim(\delta)+1\) on a star neighborhood of \(y\) in \(\mathbf X\).
This is the local-star correspondence in
\cite[Definitions~2.22--2.23 and Lemmas~2.24 and~2.35]{BMPS26}; in
particular, divisor formation commutes with passage to the local star.  The
defining PAPC condition
\cite[Definition~3.1]{BMPS26} makes its divisor by \(a\) effective.  Its
coefficient at \(\delta\) is the last member of
\eqref{eq:p053-vthree-corner-coefficient}, hence
\((\varphi\cdot z)(\tau')\geq0\).  This proves the claim for every positive
intermediate cycle \(z\), including all contractions and all nontrivial
lattice indices.  We use here only the divisor-effectivity property.  We do
not claim that \(a\circ F\) has rational additive constants in the BMPS
sense: an integral \(\mathbf R\)-affine translation may add an arbitrary real
constant, which does not enter any divisor weight.

We apply this calculation to analytic tropicalizations.  Let \(x\in U^{\an}\)
and choose an arbitrary canonical tropical chart
\((V',\operatorname{trop}_{U'})\) from the basis in
\cite[\S2.1]{GK19}, after shrinking inside \(U^{\an}\).  Restriction of the
fixed torus characters to \(U'\) induces an integral-affine map \(F\) with

\begin{equation}
 h|_{V'}=F\circ\operatorname{trop}_{U'}.
 \label{eq:p053-repaired-compatible-frame}
\end{equation}

Put \(C'=\operatorname{Trop}(U')\), with its effective analytic tropical
multiplicities, and \(\varphi=a\circ F\).  Piecewise-linear extension from a
rational subdivision gives a tropical frame for the metric of weight
\(a\circ h\).  On a common cellular refinement, \(\varphi\) is affine, hence
convex, on every maximal source cell.  This is condition~(i), in the
refined form allowed by \cite[Remark~8.5(a)]{GK19}.  Taking
\(z=[C']\) in the preceding calculation proves that every corner weight is
nonnegative, which is condition~(ii) of \cite[Theorem~8.4]{GK19}.  Only
ridges meeting the open set of the frame are involved; shrinking the open
set produces no artificial boundary weight because the criterion and all
test forms are local and compactly supported.

The compatible canonical frames just used form a basis, and the calculation
was made after an arbitrary compatible refinement.  Therefore
\cite[Theorem~8.4 and Remark~8.5(a)]{GK19} shows that the metric is psh.
Equivalently, its ordinary first Chern current is positive.  By
\cite[Definition~3.1]{GK19}, this says exactly that
\(\dprimeddprime[a\circ h]\) is a positive ordinary CLD current.  Thus
\(a\circ h\) is CLD-psh and
\eqref{eq:p053-repaired-psh-pullback} follows.  No finite-local-maximum
representation and no integral-paralinear comparison theorem is used.  No
value-group condition beyond the rational BMPS hypothesis in the statement
is required.

The argument is functorial.  Indeed, let \(f:Z\to U\) be any morphism of
algebraic varieties.  On every irreducible component and every canonical
tropical chart of \(Z^{\an}\), the composite \(h\circ f^{\an}\) is induced
by an integral-affine map of cocharacter spaces.  The preceding quotient-star
calculation applies verbatim, also when the image has smaller dimension or is
totally contracted.  The reduced components suffice because analytification
does not see nilpotents, and positivity is checked componentwise.  Hence
\(a\circ h\circ f^{\an}\) is psh for every \(f\): the metric of weight
\(a\circ h\) is functorial psh in the sense of \cite[\S3.4]{GK19}.

We now compare the delta-form and CLD products for piecewise-affine weights.
The function \(u_i=a_i\circ h\) is piecewise smooth on tropical charts.  On
every compatible cellular refinement it is affine on each cell, so its
facewise \(\mathrm d'_{\mathrm P}\mathrm d''_{\mathrm P}\)-term is zero.
Condition~(ii) of \cite[Definition~9.9]{GK17} therefore holds with ambient
superform zero, and \(\lVert1\rVert_i=\mathrm e^{-u_i}\) is a delta-metric.
The functorial psh conclusion above and \cite[Proposition~8.8]{GK19} show
that its first Chern delta-form \(\Theta_i\) is positive.

On a compatible chart as in
\eqref{eq:p053-repaired-compatible-frame}, put

\[
 \varphi_i=a_i\circ F,
 \qquad C'=\operatorname{Trop}(U'),
 \qquad c'_0=[C'],
 \qquad c'_j=\varphi_j\cdot c'_{j-1}.
\]

The stronger form of the quotient-star calculation, applied successively
to the positive intermediate cycles in any ordering, proves that every
\(c'_j\) is effective and balanced.  The same remains true after every algebraic
pullback.  The local representative of
\(\Theta_1\wedge\cdots\wedge\Theta_j\), after multiplication by the
analytic tropical integration cycle, is \(\delta_{c'_j}\).  This follows
from the construction of first Chern delta-forms in
\cite[\S\S9.13--9.14]{GK17} and commutativity of tropical divisors in
\cite[Proposition~1.14]{GK17}.  The representatives are compatible under
further refinements.  Since their arbitrary algebraic pullbacks have
effective representatives, \cite[Proposition~2.13]{GK19} shows that these
particular partial products are positive generalized delta-forms.  This
argument does not assert that arbitrary positive generalized delta-forms
have a positive wedge product.

Let \([B_j]_{\mathrm D}\) denote the ordinary CLD current underlying this
delta-form product, with \([B_0]_{\mathrm D}=[U^{\an}]\).  Inductively,
\([B_{j-1}]_{\mathrm D}\) equals the corresponding positive CLD product;
it is therefore positive and of order zero, so multiplication by the
continuous function \(u_j\) is defined.  On the compatible chart it is
represented by \(\varphi_j\delta_{c'_{j-1}}\).  Hence the weighted formula
\cite[Corollary~3.20]{GK17} gives, as an equality of ordinary currents,

\begin{equation}
 \dprimeddprime\bigl(u_j[B_{j-1}]_{\mathrm D}\bigr)
 =[B_j]_{\mathrm D}.
 \label{eq:p054-repaired-current-induction}
\end{equation}

This argument forms Chern products in the differential algebra of
delta-forms and applies the weighted tropical formula only to the associated
ordinary polyhedral current.  It does not differentiate an arbitrary element
of the generalized delta algebra.  Starting with
\([B_0]_{\mathrm D}=[U^{\an}]\), induction in
\eqref{eq:p054-repaired-current-induction} proves the first equality in
\eqref{eq:p054-repaired-delta-comparison}.  It also reconfirms positivity of
the associated ordinary current at every stage.

It remains to construct the Radon measure without applying a compact-support
result globally.  Set \(\Theta=\Theta_1\wedge\cdots\wedge\Theta_d\).  For a
relatively compact open set \(W\Subset U^{\an}\), choose a smooth compactly
supported function \(\rho\), equal to one on a neighborhood of
\(\overline W\), using \cite[Corollary~6.2.4]{CLD25}; composing with a
one-variable smooth cutoff lets us assume \(0\leq\rho\leq1\).  Then
\(\rho\Theta\in P_c^{d,d}(U^{\an})\), and
\cite[Corollary~6.15]{GK17} supplies a signed Radon measure for this compactly
supported generalized delta-form.

Restrict that measure to \(W\).  If \(\rho'\) is another such cutoff, formula
(6.15.1) of \cite{GK17} gives equality for every smooth compactly supported
test function on \(W\).  Density
\cite[Theorem~6.2.5 and Corollary~6.2.6]{CLD25} gives equality for every
continuous compactly supported test function, so the restriction is
independent of \(\rho\).  The same argument gives compatibility on overlaps;
the local measures therefore glue to a unique locally finite Radon measure
\(\mu_\delta\) on \(U^{\an}\).  Its ordinary current is the positive CLD
product already identified, so \(\mu_\delta\) is positive.  This construction
does not claim finite total mass.

For a compactly supported smooth target function \(\psi\), properness of
\(h\) makes \(\psi\circ h\) compactly supported.  Choose \(\rho\) equal to
one near this support.  On the fixed tropical chart, the compactly supported
form \(\rho\Theta\) is represented, where it is tested, by

\[
 \delta_{a_1\cdots a_d\cdot[\operatorname{Trop}(U)]}.
\]

Definition~5.8 of \cite{GK17} therefore identifies its integral against
\(\psi\circ h\) with the integral of \(\psi\) over the displayed zero-cycle.
Density of smooth compact tests in continuous compact tests proves the
pushforward identity in \eqref{eq:p054-repaired-delta-comparison}.  Notice
that the pushforward is a Radon measure because \(h\) is topologically
proper; no global finiteness is used.

Finally, let \(a_{i,k}\searrow a_i\) be arbitrary approximations as in the
statement and put \(u_{i,k}=a_{i,k}\circ h\).  The piecewise-affine case just
proved identifies \(\mu_{\delta,k}\) with the corresponding CLD mixed
product.  Proposition~8.7.2 of \cite{CLD25} therefore gives current
convergence to \(\mu_{\mathrm{CLD}}\).

We upgrade this to vague convergence.  For a compact set
\(L\subset U^{\an}\), choose a nonnegative compactly supported smooth
function \(\chi\) that is at least one on \(L\), again using a cutoff from
\cite[Corollary~6.2.4]{CLD25} followed by a smooth scalar composition.
Convergence against \(\chi\), together with positivity, gives

\[
 \sup_k\mu_{\delta,k}(L)
 \leq \sup_k\int\chi\,\mathrm d\mu_{\delta,k}<+\infty.
\]

Given a continuous compactly supported test function, approximate it
uniformly on a fixed compact neighborhood of its support by smooth compactly
supported functions.  The preceding local mass bound makes the approximation
uniform in \(k\), while current convergence handles each smooth
approximation.  This proves vague convergence
\eqref{eq:p054-repaired-source-limit} and also proves independence of the
approximating sequences.

If \(\psi\) is a continuous compactly supported function on \(\mathbf X\),
properness makes \(\psi\circ h\) compactly supported.  We may consequently
push the vague convergence through \(h\).  At every finite stage the
piecewise-affine comparison gives

\[
 h_*\mu_{\delta,k}
 =\MA_{\mathrm{poly}}(a_{1,k},\ldots,a_{d,k}),
\]

and \cite[Theorem~4.13]{BMPS26} identifies the limit of the right-hand side
with

\[
 \MA_{\mathrm{poly}}(a_1,\ldots,a_d).
\]

This proves
\eqref{eq:p054-repaired-limit-pushforward}.  The construction only supplies
a locally finite Radon-measure limit; no closure of finite delta-forms under
vague limits and no finite-total-mass assertion is used.
\end{proof}
